%% PRÄAMBEL

% -- Dokumentenklasse
\documentclass[
	% Schriftgröße
	11pt,
	% Papierformat
	a4paper,
	% Seitenverteilung
	oneside,
	% Sprache
	ngerman
	]{article}

% -- Meta-Daten
\newcommand{\mytitle}{Messtechnische Abbildung typischer Fertigungsabweichungen}
\newcommand{\myname}{Jan Doant}

% -- Seitenränder
\usepackage[left=30mm,right=20mm,top=25mm,bottom=20mm]{geometry}

% -- Zeilenabstand
\usepackage[onehalfspacing]{setspace}

% -- keine Einrückung der ersten Zeile eines neuen Absatzes 
\setlength{\parindent}{0em}

%% -- Gestaltung Kopf- und Fußzeile
% -- Kopfzeile
\usepackage{fancyhdr}
\addtolength{\headheight}{3pt}
\lhead{\mytitle}
% -- Fußzeile
\lfoot{\myname}
\cfoot{}
\rfoot{Seite~\thepage}  

%% -- deutsches Sprachpaket
\usepackage[ngerman]{babel}
% -- Umlaute
\usepackage[utf8]{inputenc}
\usepackage[T1]{fontenc}
\usepackage{ae,aecompl}
% -- deutsche Anführungszeichen
\usepackage[babel,german=quotes]{csquotes}

% -- Mathematikpakete
\usepackage{amsfonts}
\usepackage{amsmath}

% -- Grafikpaket
\usepackage{graphicx}

% -- Verweise 
% -- Verweise werden im Dokument klickbar
\usepackage{hyperref}
% -- Einfärben von Verweisen
\usepackage[dvipsnames,svgnames,x11names,hyperref]{xcolor}
\hypersetup{
	% Meta-Informationen, die in der PDF-Datei gespeichert werden
	pdfauthor={\myname},
	pdftitle={\mytitle},
	% Einfärben der Schrift
	colorlinks = true,
	% Farbe von Verknüpfungen im Text 
	linkcolor=[RGB]{63 81 181},
	% Farbe von Quellenverweisen
	citecolor= [RGB]{76 175 80},
	% Farbe von URLs
	urlcolor=[RGB]{0 76 153},
}

% -- Verweise automatisieren
%\usepackage{cleveref}
\usepackage{varioref}
\usepackage{prettyref}
% -- prettyref-Befehl konfigurieren
\newrefformat{fig}{Abbildung \vref{#1}}
\newrefformat{tab}{Tabelle \vref{#1}}

% -- Anzeigen der verwendeten Labels (nur im Dev-Modus verwenden)
\usepackage[
	% Label wird am inneren Rand angezeigt
	inner,
	% blendet Label aus, wenn final aktiv 			 
	final			 
	]{showlabels}

% -- Bibliografiestil
\usepackage[
	% Literaturverzeichnis wird in der Reihenfolge, in der die Referenzen genannt werden sortiert
	sorting=none,
	% keine Angabe der URL 		
	url=false,
	% Übersicht Biblatex-Styles:
	% https://de.sharelatex.com/learn/Biblatex_bibliography_styles 			
	style=ieee,
	% keine Angabe der ISBN 			
	isbn=false,
	% Datumsangabe nur für das Erscheiningsjahr
	date=year,
	% maximal 6 Namen in der Literaturliste
	maxbibnames=6,
	% mindestens 3 Namen in der Literaturliste    
	minbibnames=3,      
	% Verwendung von biber.exe
	backend=biber
	]{biblatex}

% -- Pfad zur .bib-Datei	
\addbibresource{bib/lit.bib} 

% -- Tabellenpaket
\usepackage{tabularx}

% -- Verzeichnisse zum Inhaltsverzeichnis hinzufügen
\usepackage[
	% Inhaltsverzeichnis selbst nicht im Inhaltsverzeichnis aufführen 
	nottoc
	]{tocbibind}

% -- TikZ
\usepackage{tikz}


%% DOKUMENT
\begin{document}

% -- Titelseite
\pagestyle{empty}
\begin{titlepage}
	
\begin{figure}[htbp]
	\begin{minipage}{0.19\textwidth}
		% \textwidth bezieht sich nun auf die Minipage
		\includegraphics[width=\textwidth]{img/logo_tuc.pdf}
	\end{minipage}
	% Auffüllen des Zwischenraums
	\hfill
	% minipage mit Text
	\begin{minipage}{0.6\textwidth} 
		\begin{center}	
			Technische Universität Chemnitz\\
			Fakultät Maschinenbau\\
			Professur Fertigungsmesstechnik
		\end{center}
	\end{minipage}
	\hfill
	\begin{minipage}{0.19\textwidth} 
		% \textwidth bezieht sich nun auf die Minipage
		\includegraphics[width=\textwidth]{img/logo_fmt.pdf}
	\end{minipage}
\end{figure}
	
	
	\vspace*{1cm}
		
		\begin{center}	
		
		\vspace{1.5cm}
		
		{\huge\textbf{Bachelorarbeit}}
		
		\vspace{1.5cm}
		
		{\large Arbeitstitel:}\\
		{\large\textbf{\mytitle}}
		
		\vspace{1.5cm}
		
		Betreuer: M. Sc. Robert Hofmann\\
		Prüfer: Univ.-Prof. Dr.-Ing. Sophie Gröger		
		
		\vfill
		
			
		
	\end{center}
	
	\begin{flushleft}
		Name: Jan Doant\\		
		E-Mail: jan.doant@s2017.tu-chemnitz.de\\
		Studiengang: Bachelor Maschinenbau (7. Semester)\\
		Matrikelnummer: 461311\linebreak
		
		
		Eingereicht am: 06.06.2018
		
	\end{flushleft}
	
	
\end{titlepage}
\newpage

% -- Selbstständigkeitserklärung
\pagestyle{empty}
\section*{Selbstständigkeitserklärung}

% -- Selbstständigkeitserklärung zum Inhaltsverzeichnis hinzufügen
\addcontentsline{toc}{section}{Selbstständigkeitserklärung}

% -- Text aus "`Hinweise zum Anfertigen wissenschaftlicher Arbeiten.pdf"' der Professur MHT
% Abschnitt 1.15
Ich erkläre, dass ich die vorliegende Arbeit selbstständig und unter Verwendung der angegebenen
Quellen und Hilfsmittel angefertigt habe. Alle Stellen, die wörtlich oder sinngemäß aus Quellen
entnommen wurden, sind als solche gekennzeichnet.

% -- Abstand der Unterschriftenleiste zum Text
\vspace{2cm}

% -- Unterschriftenleiste
\begin{tabularx}{\columnwidth}{lXr}
  \hline 
  Chemnitz, 06.06.2018 &  &Unterschrift des Verfassers
\end{tabularx}

\newpage

% -- Inhaltsverzeichnis
\pagestyle{empty}
%%Anzahl der maximal aufgeführten Gliederungspunkte
\setcounter{tocdepth}{3}
\tableofcontents
\newpage

% --Sandbox
\pagestyle{fancy}
\section{Sandbox}

\subsection{Sandbox--Zitierung}

Das ist zu eine zu zitierende Textstelle \cite{Hesse2016}.
Und hier ist noch ein Zitat.\cite{Vogt2003}
Wenn man was aus diesem Buch zitieren soll. \cite{Quarteroni2005}

\subsection{Sandbox--Grafiken einbinden}

\begin{figure}[htbp] 
  \centering
     \includegraphics[width=0.7\textwidth]{img/logo_tuc.pdf}
  \caption{Logo der Technischen Universität Chemnitz}
  \label{fig:Bild1}
\end{figure}

\newpage

\newpage


\subsection{Sandbox--Referenzierung}

Ich beziehe mich auf \prettyref{fig:Bild1}.
\newpage

% -- Einführung
\pagestyle{fancy}
\section{Einführung}
\label{sec:einfuehrung}

\subsection{Hallo}







\newpage

% -- Abbildungsverzeichnis
\pagestyle{fancy}
\listoffigures
\newpage

% -- Literaturverzeichnis
\pagestyle{fancy}
\printbibliography[
	% Literaturverzeichnis im Inhaltsverzeichnis anzeigen 
	heading=bibintoc
	]

\end{document}
